\scanpage{195}
\begin{center}- \mspage{60} -\end{center}
\margin{IX}

\uline{i.e.} \uline{définit un isomorphisme de chacun de ces schémas abéliens avec le dual de l'autre}.

La démonstration de 5.4 dans le cas général va être donnée au n° 7, comme conséquence d'une construction fort naturelle et importante due à RAYNAUD. Notons seulement ici que pour prouver 5.4, on peut supposer $S$ complet, compte tenu de 3.3.

5.5. Pour exprimer commodément les facteurs extrêmes dans les filtrations (5.1.10) et (5.1.11), dans le cas de la réduction semi-stable, introduisons les schémas en groupes constants tordus sur $S$
\[
(5.5.1)\qquad\underline M'_{\unsure{o|0}}=D(T_{\unsure{o|0}})\quad,\quad\underline M_{\unsure{o|0}}=D(T'_{\unsure{o|0}})\quad,
\]
duals des tores maximaux $T_{\unsure{o|0}},T'_{\unsure{o|0}}$ des fibres spéciales de $A,A'$ (SGA 3 X 5.7). Comme $S$ est local complet, ces schémas en groupes peuvent être considérés comme les fibres spéciales de schémas constants tordus
\[
(5.5.2)\qquad\underline M'\ ,\ \underline M\text{ sur }S\quad;
\]
ainsi, si le tore $T_{\unsure{o|0}}$ est déployé (par exemple si $k$ est séparablement clos) i.e. si $\underline M'_{\unsure{o|0}}$ est le groupe constant défini par le groupe ordinaire
\[
(5.5.3)\qquad M'=\operatorname{Hom}_{k\text{-gr}}(T_{\unsure{o|0}},\mathbb G_{mk})\quad,
\]
alors $\underline M'$ est le groupe constant défini par $M'$ :
\[
\underline M'=M'_S\quad.
\]
On peut d'ailleurs noter que, puisque $T_{\unsure{o|0}}$ et $T'_{\unsure{o|0}}$ sont isogènes (2.2.6), $T'_{\unsure{o|0}}$ est déployé si $T_{\unsure{o|0}}$ l'est, de sorte qu'il y a lieu d'introduire également
\margin{378}
\margin{378}
\margin{194}
